\SourceLeaf{253}%
\noindent die Volumengleichheit symmetrischer Körper zu beweisen. Er behauptete daraufhin, ein ähnlicher Beweis finde sich schon in \textsc{Legendres} Geometrie. Das habe ich noch nicht überprüfen können.

\Needspace{7\baselineskip}
\begin{center}
[\textsc{Gerling} \emph{an} \textsc{Gauss}. \emph{Marburg, 21. Juli 1844.}]
\end{center}

\{\mbox{.\ .\ .\ .\ .\ .\ .} Wegen der Bemerkung von Herrn \textsc{Eisenstein} über die Volumina symmetrischer Körper habe ich bei \textsc{Legendre} nachgesehen. Ich selbst habe noch die Ausgabe von 1806, hatte aber jetzt Gelegenheit, \textsc{Crelles} Übersetzung (1844) der 12. Auflage einzusehen. Gegenüber 1806 finde ich in diesen Fragen allerdings einiges geändert. --- \textsc{Legendre} hat den fraglichen Satz aber nicht \textls{unabhängig} von der Exhaustionsmethode bewiesen. Er geht nämlich von dem Satz aus, dass Pyramiden mit gleicher Grundfläche und gleicher Höhe gleich sind. --- Diesen Satz beweist er mit der Exhaustionsmethode (im Wesentlichen jetzt genau wie in meinem \textsc{Lorenz} § 330, während er sich früher enger an Euklid hielt). Dann zerlegt er die symmetrischen Polyeder in symmetrische dreiseitige Pyramiden und schließt aus jenem Satz auf ihre Gleichheit. Der Unterschied besteht also darin, dass \textsc{Legendre} die weitere Zerlegung in je 12 paarweise kongruente Pyramiden nicht ausführt. ---

Wie ich den Satz in \textsc{Lorenz} § 330 gelernt habe, weiß ich selbst nicht mehr. Selbst wenn ich ihn selbstständig gefunden haben sollte, haben mir jedenfalls Bemerkungen von \textsc{Pfaff} die Entdeckung sehr nahegelegt. Ich erinnere mich insbesondere, dass er mir (1809) vom »Segnerschen Axiom« sprach und bedauerte, dass man hier nicht tiefer eindringen könne.\}
